\scanpage{42}
On trouve donc un foncteur pl. fid
\[
 \mathcal C^\circ\lhook\joinrel\xrightarrow{\deleted{\unsure{\beta}}\alpha^\bullet}
 \mathrm{Ab}_k^{\bullet\deleted{\unsure{*}}}
\]
qui associe à tout $\deleted{\unsure{M_*}}L_*$ associe le complexe
de cochaînes $\operatorname{Hom}_{\mathcal C}(L_*,C^\bullet_*)$
\[
 0\longrightarrow\operatorname{Hom}(L_*,C^0_*)\longrightarrow
 \operatorname{Hom}(L_*,C^1_*)\longrightarrow\cdots,
\]
et dont la restriction à $\mathcal A$ est donnée par
\[
 \alpha^\bullet\deleted{\unsure{\beta}}(C^i_*)\longmapsto
 (0\to0\to\cdots\,0\to\underset{\text{degré }i}{k}
 \xrightarrow{\mathrm{id}}\underset{\text{degré }i+1}{k}\to0\to\cdots)
 =\deleted{\unsure{K(i)}}\inserted{C_i^\bullet}
\]
et par la loi
\diagram{\begin{tikzcd}[ampersand replacement=\&,column sep=1.2em,row sep=0.65em]
\& \alpha^\bullet(C^i_*) \& \alpha(C^{i+1}_*) \arrow[l,"{\deleted{\unsure{\beta}}\alpha^\bullet(d^i)}"'] \&\\
\& \Vert \& \Vert \&\\
\& \deleted{\unsure{K}}C_i^\bullet \arrow[r] \& \deleted{\unsure{K}}C_{i+1}^\bullet \& k-0\\
\& \vdots \& \vdots \&\\
\& 0 \arrow[d] \& 0 \arrow[d] \&\\
i \& k \arrow[d] \& 0 \arrow[l] \arrow[d] \&\\
i+1 \& k \arrow[d] \& k \arrow[l,"{\mathrm{id}}"'] \arrow[d] \&\\
i+2 \& 0 \arrow[d] \& k \arrow[l] \arrow[d] \&\\
\& 0 \& 0 \&\\
\& \vdots \& \vdots \&
\end{tikzcd}}
\margin{L'image essentielle \inserted{\unsure{pl.|$K^i$}} de $\mathcal C^\circ$ est la catégorie des $K^\bullet\in\operatorname{Ob}\mathrm{Ab}_k^\bullet$ qui sont libres de t. f en tte dimension, et qui sont à coh. nulle (en admettant que tt facteur direct d'un libre de t. f. est itou}
% slanted
\margin{\unsure{Commencer|Commençons|commentaire} par Ex 2}

\uline{Ex 2} Faisons $\mathcal B=\mathrm{Ab}_k^\circ$,
qui est $k$-linéaire (on suppose
$k$ commutatif), et pour tt $L_*\in\operatorname{Ob}\mathcal C$, on considère le foncteur contrav. qu'il représente, interprété
comme foncteur $k$-lin.
\[
 \mathcal C^\circ\longrightarrow\deleted{\unsure{\mathrm{Ab}^\circ_k}}\mathrm{Ab}_k^{*}
 \qquad M_*\mapsto\operatorname{Hom}_{\mathcal C}(M_*,L_*)
\]
on trouve ainsi un foncteur pl. fid (par la
\unsure{même voie|méthode|même idée|covariant}
\diagram{\begin{tikzcd}[ampersand replacement=\&,column sep=1em,row sep=1.5em]
\mathcal C \arrow[r,hook] \& \underline{\operatorname{Hom}}_{k\text{-}\mathrm{lin}}(\mathcal C^\circ,\mathrm{Ab}_k^\bullet) \arrow[r,"{\sim}"] \arrow[d,"{\simeq}"] \& \underline{\operatorname{Hom}}_{k\text{-}\mathrm{lin}}(\mathcal C,\mathrm{Ab}_k^\circ)^\circ \arrow[d,"{\simeq}"]\\
\& \underline{\operatorname{Hom}}_{k\text{-}\mathrm{lin}}(\mathcal A^\circ,\mathrm{Ab}_k^\bullet) \arrow[r,"{\sim}"] \arrow[d] \& \underline{\operatorname{Hom}}_{k\text{-}\mathrm{lin}}(\mathcal A,\mathrm{Ab}_k^\circ)^\circ \arrow[d,"{\simeq}"]\\
\& {(\mathrm{Ab}_k^\circ)_{\deleted{\unsure{*}}}^{*}\deleted{\simeq(\mathrm{Ab}_{k*}^{\unsure{*}})^\circ}} \& {(\mathrm{Ab}_k^{\circ *})^\circ}
\end{tikzcd}}
ou encore un foncteur pl. fid.
\[
 \mathcal C^{\unsure{\bullet|\circ}}\lhook\joinrel\xrightarrow{\beta_{\bullet}}\mathrm{Ab}_{k*}
\]
% archive sheet number: 41
